\documentclass[12pt]{article}
%---DOCUMENT MARGINS---
\usepackage{geometry} % Required for adjusting page dimensions and margins
\geometry{
	paper=a4paper, % Paper size, change to letterpaper for US letter size
	top=4cm, % Top margin
	bottom=2cm, % Bottom margin
	left=2cm, % Left margin
	right=2cm, % Right margin
	headheight=3cm, % Header height
	%footskip=1.5cm, % Space from the bottom margin to the baseline of the footer
	headsep=0.5cm, % Space from the top margin to the baseline of the header
	%showframe, % Uncomment to show how the type block is set on the page
}
\usepackage{cmap}

\usepackage[T1,T2A]{fontenc}
\usepackage{fontspec}
\setmainfont[
	BoldFont={* Bold},
	ItalicFont={* Italic},
	BoldItalicFont={* BoldItalic}
]{DejaVu Serif Condensed}
\setsansfont{DejaVu Sans}
\setmonofont{Dejavu Sans Mono}
\usepackage[math-style=TeX]{unicode-math}
\setmathfont{Latin Modern Math}

\usepackage[nobottomtitles*]{titlesec}
\titleformat
{\section} % command
{\fontsize{14}{14}\sffamily\bfseries} % format
{} % label
{0pt} % sep
{} % before-code
\titlespacing{\section}{0pt}{0.75em}{0.25em}
\newcommand{\tl}{}
\newcommand{\ml}{}
\newcommand{\problem}[1]{%
	\section[#1]{#1 \fontsize{14}{14}\selectfont{\hfill{\normalfont{
				\emoji{hourglass-not-done}\tl \space\space \emoji{floppy-disk} \ml}}}}
}
\AddToHook{cmd/section/before}{\clearpage}

\usepackage[dvipsnames]{xcolor}
\titleformat
{\subsection} % command
{\fontsize{14}{14}\sffamily\bfseries} % format
{\includesvg[width=0.035\textwidth, inkscapelatex=false]{lion}\space} % label
{0pt} % sep
{} % before-code
\titlespacing{\subsection}{0pt}{0.75em}{0.25em}

\setlength{\parskip}{0.5em}
\setlength{\parindent}{0pt}
\renewcommand{\baselinestretch}{1.2}
\sloppy

\usepackage{listings}
\definecolor{lstbg}{RGB}{250,250,252}
\definecolor{lstframe}{RGB}{220,220,225}
\lstdefinestyle{mystyle}{
	backgroundcolor=\color{lstbg},
	frame=single,
	rulecolor=\color{lstframe},
	basicstyle=\ttfamily\small,
	keywordstyle=\bfseries,
	breaklines=true,
	aboveskip=0.5\baselineskip,
	belowskip=0\baselineskip  
}
\lstset{style=mystyle, language=C++}

\usepackage{fancyhdr}
\pagestyle{fancy}
\setlength{\headwidth}{\textwidth}
\newcommand{\round}{}
\newcommand{\problemName}{}
\newcommand{\lang}{English}
\fancyhead[L]{
	\begin{minipage}{0.4\textwidth}
		\bf{
			EJOI 2025 \round\\
			\problemName\\
			\emoji{globe-with-meridians} \lang
		}
	\end{minipage}
	\vspace{0.5cm}
}
\fancyhead[R]{%
	\begin{minipage}{0.6\textwidth}
		\includesvg[width=\textwidth, inkscapelatex=false]{logo}
	\end{minipage}
	\vspace{0.5cm}
}
\usepackage{lastpage}
\fancyfoot[C]{\problemName\space(\lang) \thepage\ / \pageref*{LastPage}}

\raggedbottom

\usepackage{amsmath}
\usepackage{stmaryrd}

\usepackage{graphicx}
\graphicspath{{./}}
\usepackage[export]{adjustbox}
\usepackage{wrapfig}
\makeatletter
\patchcmd\WF@putfigmaybe{\lower\intextsep}{}{}{\fail}
\AddToHook{env/wrapfigure/begin}{\setlength{\intextsep}{0pt}}
\makeatother
\usepackage[inkscapearea=page, inkscapepath=./svg-inkscape]{svg}
\svgpath{{./}}

\usepackage{makecell}
\usepackage{tabularray}
\AtBeginEnvironment{table}{\vspace{-0.2cm}}
\AtEndEnvironment{table}{\vspace{-0.2cm}}
\usepackage{float}
\usepackage{placeins}
\usepackage{caption}
\captionsetup[table]{
	skip=0.25em,
	singlelinecheck=false, justification=justified
}

\usepackage{enumitem}
\setlist{itemsep=-0.4em, leftmargin=22pt, topsep=-\parskip}
\AfterEndEnvironment{itemize}{\vspace{\parskip}}
\newcommand{\tabitem}{\indent~~\llap{\textbullet}~~}

\usepackage{hyperref}
\hypersetup{
	colorlinks=true,
	citecolor=blue,
	linkcolor=blue,
	urlcolor=cyan,
}

\usepackage{emoji}

\usepackage[english]{babel}

\begin{document}

\renewcommand{\round}{Day 1}
\renewcommand{\problemName}{Task Prison}
\renewcommand{\lang}{Armenian}

\renewcommand{\tl}{$2$ sec.}
\renewcommand{\ml}{$1024$ MB}
\problem{\problemName}

Ալիսն ու Բոբը անարդար կերպով դատապարտվել են շատ խիստ պահվող  բանտում։ Այժմ նրանք պետք է պլանավորեն իրենց փախուստը։ Դրա համար նրանք պետք է կարողանան հաղորդակցվել հնարավորինս արդյունավետ (մասնավորապես, Ալիսը պետք է ամեն օր տեղեկություններ ուղարկի Բոբին)։ Բայց նրանք չեն կարող հանդիպել և միայն կարող են տեղեկություն փոխանակել անձեռոցիկների վրա գրավոր նշումների միջոցով։ Ամեն օր Ալիսը ուզում է ուղարկել նոր տեղեկություն Բոբին՝ $0$ - ի և $N-1$ - ի միջև թիվ։ Ամեն ճաշի ժամանակ Ալիսը ստանում է երեք անձեռոցիկներ և յուրաքանչյուր անձեռոցիկի վրա գրում է $0$ - ի և $M-1$ - ի միջև մեկ թիվ (կարող են լինել կրկնություններ) և թողնում դրանք իր նստարանին։ Այնուհետև, նրանց թշնամին՝ Չարլին, ոչնչացնում է մեկ անձեռոցիկը և խառնում մնացած երկուսը։ Վերջում Բոբը գտնում է երկու մնացած անձեռոցիկները և կարդում դրանց վրա գրված թվերը։ Նա պետք է ճշգրիտ պարզի անհրաժեշտ տեղեկատվությունը, որը Ալիսը ցանկանում էր ուղարկել նրան։ Անձեռոցիկների վրա տեղը սահմանափակ է, ուստի $M$-ը ֆիքսված է։ Բայց Ալիսի ու Բոբի նպատակն է մաքսիմիզացնել ուղարկվող տեղեկատվությունը, այնպես որ նրանք ազատ են ընտրելու $N$-ը որքան կարող են մեծ։ Օգնեք Ալիսին և Բոբին իրականացնելու այնպիսի ռազմավարություն յուրաքանչյուրի համար, որպեսզի $N$-ի արժեքը հնարավորինս մեծ լինի։

\subsection{Իրականացման մանրամասներ}

Քանի որ սա հաղորդակցման խնդիր է, ձեր ծրագիրը կաշխատի երկու առանձին գործարկումներով (մեկը Ալիսի համար և մեկը Բոբի համար), որոնք չեն կարող ունենալ ընդհանուր տվյալներ կամ հաղորդակցվել որևէ այլ կերպ, բացի այստեղ նկարագրվածից: Դուք պետք է իրականացնեք երեք ֆունկցիա․
\begin{lstlisting}
 int setup(int M);
\end{lstlisting}

Սա կանչվելու է մեկ անգամ, նախքան ձեր ծրագիրը կաշխատացվի Ալիսի համար և մեկ անգամ, նախքան ձեր ծրագիրը կաշխատացվի Բոբի համար։ Այս ֆունկցիային տրվում է $M$-ը և ինքը պետք է վերադարձնի ցանկալի $N$-ը։ \texttt{setup}-ի երկու կանչերն էլ պետք է վերադարձնեն միևնույն  $N$-ը։

\begin{lstlisting}
 std::vector<int> encode(int A);
\end{lstlisting}

Սա իրականացնում է Ալիսի ռազմավարությունը։ Այն պետք է կոդավորի  $A$ ($0 \leq A < N$) թիվը և պետք է վերադարձնի երեք թիվ՝ $W_1, W_2, W_3$ ($0 \leq W_i < M$), որոնք կոդավորում են $A$-ն։ Այս ֆունկցիան կանչվելու է ընադամենը $T$ անգամ՝ ամեն օր մեկ անգամ ($A$-ի արժեքները ինչ որ օրերում կարող են կրկնվել)։

\begin{lstlisting}
 int decode(int X, int Y);
\end{lstlisting}

Սա իրականացնում է Բոբի ռազմավարությունը։ Այն կանչվելու է  \texttt{encode}-ի վերադարձրած թվերից երկուսի վրա, որոնք տրվելու են կամայական հերթականությամբ։ Այն պետք է վերադարձնի նույն $A$ արժեքը, որը \texttt{encode}-ը ստացել է։ Այս ֆունկցիան նույնպես կանչվելու է $T$ անգամ՝ \texttt{encode}-ի $T$ կանչերին համապատասխան, նրանք լինելու են նույն հերթականությամբ։ \texttt{encode}-ի բոլոր կանչերը նախորդելու են  \texttt{decode}-ի բոլոր կանչերին։

\subsection{Սահմանափակումներ}

\begin{itemize}
\item $M \leq 4300$
\item $T = 5000$
\end{itemize}

\subsection{Գնահատումը}

Յուրաքանչյուր ենթախնդրի համար, դուք կստանաք որպես միավոր մի S կոտորակ, որը կախված է \texttt{setup}-ի կողմից այդ ենթախնդրի բոլոր թեստերի համար վերադարձված փոքրագույն $N$-ից: Այն նաև կախված է $N^*$-ից, որը $N$-ի նպատակային արժեքն է, որի դեպքում տրվում է տվյալ ենթախնդրի համար նախատեսված միավորն ամբողջությամբ:

\begin{itemize}
\item Եթե ձեր ծրագիրը որևէ թեստի համար սխալ է աշխատում, ապա  $S = 0$։
\item Եթե $N \geq N^*$, ապա $S = 1.0$։
\item Եթե $N < N^*$, ապա $S = \max\left(0.35 \max\left(\frac{\log(N) - 0.985 \log(M)}{\log(N^*) - 0.985 \log(M)}, 0.0\right)^{0.3} + 0.65 \left(\frac{N}{N^*}\right)^{2.4}, 0.01\right)$։
\end{itemize}

\subsection{Ենթախնդիրներ}

\begin{table}[H]
\begin{tblr}{|Q[c,m]|Q[c,m]|Q[c,m]|Q[c,m]|}
\hline
\textbf{Ենթախնդիր} & \textbf{Միավոր} & $M$ & $N^*$ \\
\hline
$1$ & $10$ & $700$ & $82017$ \\
\hline
$2$ & $10$ & $1100$ & $202217$ \\
\hline
$3$ & $10$ & $1500$ & $375751$ \\
\hline
$4$ & $10$ & $1900$ & $602617$ \\
\hline
$5$ & $10$ & $2300$ & $882817$ \\
\hline
$6$ & $10$ & $2700$ & $1216351$ \\
\hline
$7$ & $10$ & $3100$ & $1603217$ \\
\hline
$8$ & $10$ & $3500$ & $2043417$ \\
\hline
$9$ & $10$ & $3900$ & $2536951$ \\
\hline
$10$ & $10$ & $4300$ & $3083817$ \\
\hline
\end{tblr}
\end{table}
\FloatBarrier

\subsection{Օրինակ}

Դիտարկենք հետևյալ օրինակը, որտեղ $T = 5$։ Այստեղ մենք ունենք կոդավորման այսպիսի սխեմա․ Եթե Ալիսն ուղարկում է երեք հատ նույն թիվը, ապա նա կոդավորում է $0$ թիվը, իսկ եթե երեք տարբեր թվեր, ապա նա կոդավորում է $1$ թիվը։ Նկատենք, որ Բոբը կարող է ապակոդավորել սկզբնական թիվը Ալիսի ուղարկած երեք թվերից ցանկացած երկուսի միջոցով։

\begin{table}[H]
\begin{tblr}{|Q[c,m]|Q[c,m]|Q[c,m]|}
\hline
\textbf{\makecell{Execution}} & \textbf{\makecell{Function call}} & \textbf{Return value} \\
\hline
Alice & \texttt{setup(10)} & \texttt{2} \\
\hline
Bob & \texttt{setup(10)} & \texttt{2} \\
\hline
Alice & \texttt{encode(0)} & \texttt{\{5, 5, 5\}} \\
\hline
Alice & \texttt{encode(1)} & \texttt{\{8, 3, 7\}} \\
\hline
Alice & \texttt{encode(1)} & \texttt{\{0, 3, 1\}} \\
\hline
Alice & \texttt{encode(0)} & \texttt{\{7, 7, 7\}} \\
\hline
Alice & \texttt{encode(1)} & \texttt{\{6, 2, 0\}} \\
\hline
Bob & \texttt{decode(5, 5)} & \texttt{0} \\
\hline
Bob & \texttt{decode(8, 7)} & \texttt{1} \\
\hline
Bob & \texttt{decode(3, 0)} & \texttt{1} \\
\hline
Bob & \texttt{decode(7, 7)} & \texttt{0} \\
\hline
Bob & \texttt{decode(2, 0)} & \texttt{1} \\
\hline
\end{tblr}
\end{table}

\subsection{Գրեյդերի նմուշը}

Գրեյդերի նմուշի համար \texttt{encode}-ի և \texttt{decode}-ի բոլոր կանչերը լինելու են ձեր ծրագրի միևնույն կատարման ժամանակ։ Նաև  \texttt{setup}-ը կանչվելու է միայն մեկ անգամ (հակառակ իսկական գրեյդերի, որի դեպքում յուրաքանչյուր կատարման համար մեկ անգամ է կանչվում)։

Մուտքում տրվում է միայն մեկ թիվ՝  $M$-ը։ Ապա այն կտպի ձեր   \texttt{setup}-ի վերադարձրած  $N$-ը։ Հետո այն կանչելու է  \texttt{encode} և \texttt{decode} ֆունկցիաները այս հերթականությամբ  $T$ անգամ պատահական գեներացված $0$-ից $N-1$ սահմաններում գտնվող թվերով, նաև պատահական է որոշվում, թե  \texttt{encode}-ի տված երեք թվերից որ երկուսը և ինչ հերթականությամբ տրվեն \texttt{decode}-ին։ Այն կտպի հաղորդագրություն սխալի մասին, եթե ձեր ծրագիրը սխալ աշխատի։

\end{document}
